\documentclass[a4paper]{article}
\usepackage{amsfonts}
\usepackage{amsmath}
\usepackage{amssymb}
\usepackage{graphicx}     

\begin{document}
  
\noindent \large Auteur: Thomas Eertmans \\
\noindent \large Studierichting: Informatica\\
\noindent \large Oefeningenreeks 1C nr. 31.5\\
\newcommand{\Bgtan}{\operatorname{Bgtan}}
\newcommand{\cis}{\operatorname{cis}}
\medskip

\normalsize

Bereken in $\mathbb{C}$ door gebruik van goniometrische vormen: $\left(\sqrt{6}+i\sqrt{2}\right)^{12}$ \\

$r = \sqrt{8}$ \\

$\tan \theta = \dfrac{\sqrt{2}}{\sqrt{6}}$ \\

$\theta = \Bgtan \dfrac{\sqrt{3}}{3} = \dfrac{\pi}{6}$ \\

$\left(\sqrt{8} \cdot \cis \left(\dfrac{\pi}{6}\right)\right)^{12}$ \\

$= \sqrt{8}^{12} \cis \left(12 \cdot \dfrac{\pi}{6}\right)$ \\

$= 8^6 \cdot \cis ( 2\pi)$ \\

$= 8^6 \cdot \cis (0)$ \\

$= 8^6 \cdot 1 + 8^6i \cdot 0$ \\

$= 262144$ \\

\begin{thebibliography}{999}
%\bibitem{a} Referentie
\end{thebibliography}
\end{document} 
